% 本文件是实习报告的附录：核心代码清单
% 论文的主文件位于上级目录的 `main.tex`

\section{附录：核心代码清单} \label{app:code}

为控制篇幅（核心代码不超过 3 页），本附录仅给出两段最具代表性的核心代码：PCA 法向量估计，
以及“由粗到细”分层细分的 \texttt{segment} 函数（区域生长 \texttt{regionGrow} 已在正文
\autoref{lst:rg} 给出）。KD 树、文件读写、图形界面等完整源码见随附仓库，结构与构建方法见
仓库根目录 \texttt{README.md}。

\subsection{PCA 法向量与曲率估计}
\lstinputlisting[language=C++, caption={Normals.cpp —— PCA 估计法向量与曲率}]{code/Normals.cpp}

\subsection{由粗到细分层细分}
\lstinputlisting[language=C++, caption={Segmentation.cpp —— segment 函数（由粗到细）},
  firstline=90, lastline=144]{code/Segmentation.cpp}
